\begin{abstract}
Embedding text in one language within text of another is commonplace
for numerous purposes,
but usually requires tedious and error-prone
``escaping'' transformations on the embedded string.
We propose a simple cross-language syntactic discipline,
\emph{matchertext},
embodied in the single rule that \emph{matchers must match} --
where \emph{matchers} are the ASCII character pairs
\verb|()|, \verb|[]|, and \verb|{}|.
This discipline
enables the safe embedding any matchertext-compliant string
into any matchertext-aware language via simple ``copy-and-paste'' --
with no escaping, obfuscation, or expansion of embedded strings.
We formalize this guarantee in a machine-checked Lean proof
that delimiting valid matchertext always yields valid matchertext,
so embeddings compose to any depth.
Simple string-literal syntax extensions like \verb|M"(|$\dots$\verb|)"|,
backward-compatible with a wide variety of languages,
could enable the robust embedding of arbitrary matchertext strings
consistently without escaping.
We apply this syntactic discipline experimentally
to several common and frequently-embedded
language syntaxes such as URIs, HTML, and JavaScript,
exploring the benefits, costs, and compatibility issues
in adopting the proposed matchertext discipline.
Because breaking out of a matcher pair would require an unmatched matcher,
delimiting untrusted input with matchers rather than quotes
gives a uniform, host-independent structural defense
against injection attacks such as SQL injection and cross-site scripting.
An empirical analysis of roughly 470~million string literals
drawn from thousands of production repositories
finds that real-world code is already overwhelmingly compliant
with this discipline,
and that adopting it could eliminate
around 95\% of the backslash escapes in string literals.
We also present MinML,
an experimental markup syntax designed around matchertext,
as a concise but general alternative syntax for writing HTML or XML.
\end{abstract}
